\documentclass{beamer}

% Theme choice
\usetheme{Madrid} % You can change to e.g., Warsaw, Berlin, CambridgeUS, etc.

% Encoding and font
\usepackage[utf8]{inputenc}
\usepackage[T1]{fontenc}

% Graphics and tables
\usepackage{graphicx}
\usepackage{booktabs}

% Code listings
\usepackage{listings}
\lstset{
    language=Python, % Set default language to Python
    basicstyle=\ttfamily\small,
    keywordstyle=\color{blue}\bfseries,
    commentstyle=\color{gray}\itshape,
    stringstyle=\color{red},
    breaklines=true,
    frame=tb, % top and bottom frame lines
    showstringspaces=false,
    numbers=left,
    numberstyle=\tiny\color{gray},
    tabsize=4,
    captionpos=b % caption position at the bottom
}

% Math packages
\usepackage{amsmath}
\usepackage{amssymb}

% Colors
\usepackage{xcolor}

% TikZ and PGFPlots
\usepackage{tikz}
\usepackage{pgfplots}
\pgfplotsset{compat=1.18}
\usetikzlibrary{positioning}

% Hyperlinks
\usepackage{hyperref}

% Title information
\title{Chapter 7: Functions: Creating Reusable Code}
\author{Teaching Assistant}
\institute{University Department}
\date{\today}

\begin{document}

\frame{\titlepage}

\begin{frame}[fragile]
    \frametitle{Chapter 7: Functions: Creating Reusable Code}
    \framesubtitle{The Power of Reusability}

    \begin{itemize}
        \item Welcome to one of the most fundamental concepts in programming: \textbf{functions}.
        \item Imagine building a complex machine. Instead of creating every single screw and gear from scratch each time, you would create a blueprint and manufacture them as needed.
    \end{itemize}
    \vspace{1cm}

    \begin{block}{The Core Analogy}
        Functions are the "blueprints" of your code. They allow you to define a piece of logic once and reuse it whenever necessary.
    \end{block>

\end{frame}

\begin{frame}[fragile]
    \frametitle{What is a Function? \& Chapter Goals}

    \begin{itemize}
        \item A \textbf{function} is a named, self-contained block of code that performs a specific task.
        \item You write the code once inside the function, and then "call" it anytime you need to perform that task.
    \end{itemize}

    \begin{block}{Key Principle: DRY (Don't Repeat Yourself)}
        \centering
        \textbf{Write Once, Use Many Times.}
        \vspace{0.2cm}
        
        This is the key to managing complexity and building powerful, maintainable programs.
    \end{block}
    \vspace{0.5cm}

    \textbf{In this chapter, you will learn:}
    \begin{itemize}
        \item \textbf{Defining and Calling Functions:} The syntax for creating and executing functions.
        \item \textbf{Parameters and Arguments:} How to pass information \textit{into} a function.
        \item \textbf{Return Values:} How to get results \textit{out of} a function.
        \item \textbf{Scope:} Where variables are accessible (local vs. global environments).
    \end{itemize}
\end{frame}

\begin{frame}[fragile]
    \frametitle{Why Do We Need Functions?}
    \begin{block}{The Core Problem: Repetitive Code}
        Imagine you copy and paste the same 10 lines of code in five different places. If you need to fix a bug or make an update, you must find and change all five copies perfectly. This is inefficient and leads to errors.
    \end{block}

    \begin{block}{The Solution: The DRY Principle}
        Functions allow us to follow a core principle of software development: \textbf{DRY (Don't Repeat Yourself)}. You write the code once, inside a function, and then call it whenever you need it.
    \end{block}

    \vspace{1em}
    Functions provide four main benefits:
    \begin{itemize}
        \item Reduces Redundancy
        \item Improves Readability
        \item Simplifies Debugging
        \item Enables Abstraction
    \end{itemize}
\end{frame}

\begin{frame}[fragile]
    \frametitle{Benefits: Redundancy \& Readability}
    \begin{columns}[T] % The T option aligns the columns at the top
        \begin{column}{0.5\textwidth}
            \begin{block}{1. Reduces Redundancy}
                Define logic once and reuse it. If you need to update it (e.g., change the value of pi), you only have to make one change.
                \vfill
                \textbf{Example: With a Function}
\begin{verbatim}
def calculate_circle_area(radius):
    return 3.14159 * radius**2

# Call the same logic multiple times
area1 = calculate_circle_area(5)
area2 = calculate_circle_area(10)
\end{verbatim}
            \end{block}
        \end{column}
        \begin{column}{0.5\textwidth}
            \begin{block}{2. Improves Readability}
                Well-named functions make your code read like a summary of steps, hiding the complex details.
                \vfill
                \textbf{Example: Organized Main Logic}
\begin{verbatim}
def display_welcome_message():
    #... implementation details ...

def get_user_data():
    #... implementation details ...

def calculate_birth_year(age):
    #... implementation details ...

# Main code is clean and readable
display_welcome_message()
name, age = get_user_data()
year = calculate_birth_year(age)
\end{verbatim}
            \end{block}
        \end{column}
    \end{columns}
\end{frame}

\begin{frame}[fragile]
    \frametitle{Benefits: Debugging \& Abstraction}
    \begin{columns}[T]
        \begin{column}{0.5\textwidth}
            \begin{block}{3. Simplifies Debugging}
                Functions isolate code into small, independent units.
                \begin{itemize}
                    \item It's much easier to find a bug in a 10-line function than in a 500-line script.
                    \item Error messages often point directly to the function that failed, narrowing your search.
                    \item If you know the input is correct but the output is wrong, the bug \textit{must} be inside that specific function.
                \end{itemize}
            \end{block}
        \end{column}
        \begin{column}{0.5\textwidth}
            \begin{block}{4. Enables Abstraction}
                Abstraction means hiding complexity and showing only what's necessary. Functions are like ``black boxes''.
                \vfill
                \textbf{Analogy: Driving a Car}\\
                You use the steering wheel and pedals (the interface) without needing to know how the engine works (the implementation).
                \vfill
                \textbf{Example: Python's \texttt{len()} function}}
\begin{verbatim}
# We use len() without seeing its code.
# This is abstraction in action.
numbers = [10, 20, 30, 40]
count = len(numbers)
print(count) # Output: 4
\end{verbatim}
            \end{block}
        \end{column}
    \end{columns}
\end{frame}

\begin{frame}[fragile]
    \frametitle{Anatomy of a Python Function - The Blueprint}
    \framesubtitle{Understanding the Structure}
    
    A function is like a recipe for a task. It consists of a \textbf{header} (the first line) and an indented \textbf{body} (the instructions).
    
    \begin{block}{Example: `calculate_sum`}
\begin{lstlisting}[language=Python]
# The Function Header
def calculate_sum(num1, num2):
    # The Function Body (indented)
    total = num1 + num2
    return total
\end{lstlisting}
    \end{block}
    
    \begin{itemize}
        \item The \textbf{Header} defines the function's name and its parameters.
        \item The \textbf{Body} contains the code that runs when the function is called.
    \end{itemize}
    
\end{frame}

\begin{frame}[fragile]
    \frametitle{Anatomy of a Python Function - The Header}
    \framesubtitle{\texttt{def calculate\_sum(num1, num2):}}
    
    The header is the signature of the function. Let's break it down:
    \begin{enumerate}
        \item \textbf{\texttt{def} Keyword:} Stands for "define". It tells Python you are creating a new function.
        
        \item \textbf{Function Name (\texttt{calculate\_sum}):} A descriptive name following the \texttt{snake\_case} convention (lowercase with underscores).
        
        \item \textbf{Parentheses \texttt{()}:} Required after the name. They hold the function's parameters. They must be included even if there are no parameters (\texttt{def my\_function():}).
        
        \item \textbf{Parameters (\texttt{num1, num2}):} Placeholder variables for the inputs (arguments) the function will receive.
        
        \item \textbf{Colon \texttt{:}:} Marks the end of the header and signals that an indented code block (the body) will follow.
    \end{enumerate}
    
\end{frame}

\begin{frame}[fragile]
    \frametitle{Anatomy of a Python Function - The Body}
    \framesubtitle{The Indented Instructions}
    
    The body contains the logic that executes when the function is called.
    
    \begin{itemize}
        \item \textbf{Indented Body:}
        \begin{itemize}
            \item The block of code that forms the function's logic.
            \item \textbf{Crucially, it must be indented} (standard is 4 spaces).
            \item Indentation is how Python groups statements together.
        \end{itemize}
        \bigskip % Adds a bit of vertical space
        
        \item \textbf{\texttt{return} Statement (Optional):}
        \begin{itemize}
            \item Sends a value back from the function to the calling code.
            \item Immediately exits the function once executed.
            \item If omitted, the function automatically returns the special value \texttt{None}.
        \end{itemize}
    \end{itemize}
    
\end{frame}

\begin{frame}[fragile]
    \frametitle{Defining and Calling a Function - Part 1: The Concept}

    A function is a stored set of instructions. Using a function is a two-step process:
    \begin{itemize}
        \item \textbf{Defining} a function is like writing down a recipe; it creates the instructions but doesn't do anything yet.
        \item \textbf{Calling} a function is like following the recipe to actually cook the dish; it executes the instructions.
    \end{itemize}

    \begin{block}{Step 1: Define the Function}
        The \texttt{def} keyword tells Python you are creating a function. Python reads the indented code and stores it in memory. \textbf{The code does not run at this stage.}
    \end{block}
    
    \begin{lstlisting}[language=Python]
# This code is read and stored under the name 'display_greeting'.
# The print statements inside are NOT executed yet.
def display_greeting():
    print("---------------------")
    print("Welcome to the program!")
    print("---------------------")

# Program execution continues here...
# (Nothing has been printed to the screen yet)
    \end{lstlisting}

\end{frame}

\begin{frame}[fragile]
    \frametitle{Defining and Calling a Function - Part 2: The Execution}
    
    \begin{block}{Step 2: Call the Function}
        To execute the stored code, you must \textbf{call} it by writing its name followed by parentheses \texttt{()}. This is the trigger!
    \end{block}

    \begin{columns}[T]
        \begin{column}{0.5\textwidth}
            \textbf{Full Code:}
            \begin{lstlisting}[language=Python]
# 1. DEFINE the function
def display_greeting():
    print("---------------------")
    print("Welcome to the program!")
    print("---------------------")

# Other code could run here...

# 2. CALL the function
# Execution jumps into the
# function's code block NOW.
display_greeting()
            \end{lstlisting}
        \end{column}
        \begin{column}{0.5\textwidth}
            \textbf{Program Output:}
            \begin{verbatim}
---------------------
Welcome to the program!
---------------------
            \end{verbatim}
            
            \begin{alertblock}{The Power of Reusability}
            Define the function once, and you can call it as many times as you need!
            \begin{lstlisting}[language=Python]
display_greeting() # First call
# ... do other work ...
display_greeting() # Second call
            \end{lstlisting}
            \end{alertblock}
        \end{column}
    \end{columns}
\end{frame}

\begin{frame}[fragile]
    \frametitle{Passing Information with Parameters}
    
    Parameters make functions flexible and reusable. They act as placeholders for data that a function needs to do its work.
    
    \begin{itemize}
        \item \textbf{Parameter}: A variable listed inside the parentheses in the function definition. It's a placeholder.
        \item \textbf{Argument}: The actual value that is sent to the function when it is called.
    \end{itemize}
    
    \begin{block}{Analogy: The Vending Machine}
        \begin{itemize}
            \item \textbf{Function}: The vending machine itself (\texttt{get\_snack}).
            \item \textbf{Parameter}: The placeholder for your choice (e.g., \texttt{item\_code}).
            \item \textbf{Argument}: The specific choice you make (e.g., the string \texttt{"B4"}).
        \end{itemize}
    \end{block}
    
\end{frame}

\begin{frame}[fragile]
    \frametitle{Passing Information with Parameters - Code Example}
    
    \begin{lstlisting}[language=Python]
# 'name' is the PARAMETER (a placeholder in the function's definition)
def personalized_greet(name):
    # Inside the function, 'name' behaves like a variable
    # holding the value that was passed in.
    print(f"Hello, {name}! Welcome.")

# "Alice" is the ARGUMENT (the actual data passed during the call)
personalized_greet("Alice")
# Python effectively does: name = "Alice"
# Output: Hello, Alice! Welcome.

# "Bob" is another ARGUMENT for a separate function call
personalized_greet("Bob")
# Python effectively does: name = "Bob"
# Output: Hello, Bob! Welcome.
    \end{lstlisting}
    
    \begin{block}{Key Takeaway}
        By using a parameter, we created one \texttt{personalized\_greet} function that can greet anyone, making our code highly reusable.
    \end{block}
    
\end{frame}

\begin{frame}[fragile]
    \frametitle{Getting Results with the \texttt{return} Statement}
    
    Functions can perform a task and send a result back to the code that called it using the \texttt{return} keyword.
    
    \begin{block}{Capturing the Output}
        The returned value can be stored in a variable, printed, or used in another expression.
    \end{block}
    
    \begin{lstlisting}[language=Python]
# This function calculates and returns a value
def calculate_area(length, width):
    area = length * width
    return area

# Call the function and store its output in a variable
room_area = calculate_area(10, 5) # room_area is now 50
print(f"The area of the room is {room_area} sq ft.")
    \end{lstlisting}
    
\end{frame}

\begin{frame}[fragile]
    \frametitle{Getting Results with the \texttt{return} Statement - Data Flow}
    
    Let's trace the flow of data from the function call to the final result.
    
    \begin{lstlisting}[language=Python]
# This function calculates and returns a value
def calculate_area(length, width):
    area = length * width
    return area  # Step 2: The value of 'area' (50) is sent back

# Step 1: Call the function. Python pauses here and runs the function.
# Step 3: The function call is replaced by its returned value (50).
#         The line effectively becomes: room_area = 50
room_area = calculate_area(10, 5)

# Step 4: Use the captured value.
print(f"The area of the room is {room_area} sq ft.")
# Output: The area of the room is 50 sq ft.
    \end{lstlisting}
    
    \begin{block}{Key Rules for \texttt{return}}
        \begin{itemize}
            \item \textbf{Immediate Exit:} When a \texttt{return} statement is executed, the function stops immediately. No code after it will run.
            \item \textbf{Implicit Return:} A function that ends without an explicit \texttt{return} automatically returns a special value: \texttt{None}.
        \end{itemize}
    \end{block}
    
\end{frame}

\begin{frame}[fragile]
    \frametitle{Understanding Variable Scope: The Core Idea}
    \textbf{Scope} defines where a variable can be seen and accessed in your code.
    
    \begin{block}{An Analogy: The House and The Rooms}
        Think of your entire program as a large \textbf{house}. Each function you define is like a separate, private \textbf{room} inside that house.
    \end{block}
    
    There are two primary types of scope:
    \begin{itemize}
        \item \textbf{Global Scope (The House):} Variables created here are accessible from anywhere in the program, including inside any "room".
        
        \item \textbf{Local Scope (A Room):} Variables created inside a function are private to that "room". They cannot be seen or used from the main house or other rooms.
    \end{itemize}
    
\end{frame}

\begin{frame}[fragile]
    \frametitle{Variable Scope: Local vs. Global in Code}
    \begin{itemize}
    \item \textbf{Local Scope:} Variables created \textit{inside} a function are local. They only exist while the function is running.
    \item \textbf{Global Scope:} Variables created \textit{outside} of any function are global. They can be accessed from anywhere in the script.
    \end{itemize}
    
    \begin{block}{Example}
    \begin{lstlisting}[language=Python]
global_var = "I am global"

def my_function():
    # This variable is created in the local scope
    local_var = "I am local" 
    
    # We can access global variables from inside a function
    print(global_var)

my_function()

# This line would cause an error because local_var
# does not exist in the global scope.
# print(local_var)
    \end{lstlisting}
    \end{block}
\end{frame}

\begin{frame}[fragile]
    \frametitle{Scope: The \texttt{NameError} and Best Practices}
    
    \begin{block}{Why does \texttt{print(local\_var)} fail?}
        \begin{itemize}
            \item The variable \texttt{local\_var} is created inside \texttt{my\_function()}. It lives only in the function's local scope (its "room").
            \item When the function finishes, its local scope is destroyed, and \texttt{local\_var} vanishes completely.
            \item Trying to access it from the global scope results in a \textbf{\texttt{NameError}}, as Python cannot find a variable with that name there.
        \end{itemize}
    \end{block}
    
    \textbf{Best Practices:}
    \begin{enumerate}
        \item \textbf{Encapsulation is Key:} Local scopes keep functions self-contained and prevent them from accidentally interfering with each other.
        \item \textbf{Pass Data Explicitly:}
        \begin{itemize}
            \item Pass data \textbf{INTO} functions using \textbf{parameters}.
            \item Get data \textbf{OUT OF} functions using the \textbf{\texttt{return}} statement.
        \end{itemize}
        \item \textbf{Use Globals Sparingly:} Relying on global variables can make your code harder to read, debug, and maintain.
    \end{enumerate}
    
\end{frame}

\begin{frame}
    \frametitle{The Power of Decomposition (Part 1) - The Concept}
    
    Problem decomposition is the process of breaking a complex problem down into smaller, simpler, and more manageable sub-problems. Each sub-problem can then be solved with a dedicated function.
    
    \begin{block}{Example Problem}
        Calculate the final price of an item after a 10\% discount and an 8\% sales tax are applied.
    \end{block}
    
    \textbf{Decomposition into Steps:}
    \begin{enumerate}
        \item A function to calculate the price after the discount is applied.\\
        \texttt{calculate\_discounted\_price(price, percent)}
        
        \vspace{0.5cm}
        
        \item A function to calculate the final price by adding tax to the discounted price.\\
        \texttt{add\_tax(price, tax\_rate)}
    \end{enumerate}
    
\end{frame}

\begin{frame}[fragile]
    \frametitle{The Power of Decomposition (Part 2) - Code Example}
    
    The main script acts as an "orchestrator," coordinating the work of specialized functions.
    
    \begin{lstlisting}[language=Python]
# Function for sub-problem 1: Applying the discount
def calculate_discounted_price(price, percent):
    discount = price * (percent / 100)
    return price - discount

# Function for sub-problem 2: Adding the tax
def add_tax(price, tax_rate):
    tax = price * (tax_rate / 100)
    return price + tax

# --- Main Code: The Orchestrator ---
original_price = 100.0

# Step 1: Call the discount function
price_after_discount = calculate_discounted_price(original_price, 10)

# Step 2: Use the result to call the tax function
final_price = add_tax(price_after_discount, 8)

print(f"Final Price: ${final_price:.2f}") # $97.20
    \end{lstlisting}
\end{frame}

\begin{frame}
    \frametitle{The Power of Decomposition (Part 3) - Key Benefits}
    
    Decomposing problems into functions makes your code better in several ways:
    
    \begin{itemize}
        \item \textbf{Clarity \& Readability}
        \begin{itemize}
            \item The main code reads like a high-level summary of steps.
            \item It focuses on \textit{what} is happening, not \textit{how}. The complex details are hidden inside functions.
        \end{itemize}
        
        \vspace{0.5cm}
        
        \item \textbf{Reusability}
        \begin{itemize}
            \item Functions like \texttt{add\_tax} are reusable tools.
            \item You can call them from anywhere in your program without rewriting the logic.
        \end{itemize}
        
        \vspace{0.5cm}
        
        \item \textbf{Easier Debugging}
        \begin{itemize}
            \item If the tax calculation is wrong, you know exactly which small function to examine (\texttt{add\_tax}).
            \item This isolates bugs and makes them much faster to find and fix.
        \end{itemize}
    \end{itemize}
    
\end{frame}

\begin{frame}[fragile]
    \frametitle{Putting It All Together: A Complete Example}
    
    Let's combine our knowledge of functions to build a small, organized program that converts temperature from Celsius to Fahrenheit. This example demonstrates how to structure a program by decomposing a problem.
    
    \begin{block}{Temperature Conversion Program}
\begin{lstlisting}[language=Python]
def get_celsius_input():
    """Gets temperature in Celsius from the user."""
    c_temp = float(input("Enter temp in Celsius: "))
    return c_temp

def convert_c_to_f(celsius_temp):
    """Converts Celsius to Fahrenheit."""
    fahrenheit = (celsius_temp * 9/5) + 32
    return fahrenheit

def display_result(celsius, fahrenheit):
    """Prints the final result."""
    print(f"{celsius:.1f}C is equal to {fahrenheit:.1f}F.")

# --- Main part of the script ---
celsius = get_celsius_input()
fahrenheit = convert_c_to_f(celsius)
display_result(celsius, fahrenheit)
\end{lstlisting}
    \end{block}
\end{frame}

\begin{frame}[fragile]
    \frametitle{Putting It All Together: Code Analysis}
    
    This program demonstrates \textbf{problem decomposition}. We've broken the task into three logical parts, each handled by a dedicated function:
    
    \begin{enumerate}
        \item \textbf{Get Input:} Get the temperature from the user.
        \item \textbf{Process Data:} Perform the conversion calculation.
        \item \textbf{Display Output:} Show the final result to the user.
    \end{enumerate}
    
    \begin{block}{Function Responsibilities}
    \begin{itemize}
        \item \texttt{get\_celsius\_input()}:
            \begin{itemize}
                \item \textbf{Responsibility:} Manages all user interaction for input.
                \item \textbf{Action:} Prompts user, converts input to a \texttt{float}, and returns the value.
            \end{itemize}
        \item \texttt{convert\_c\_to\_f()}:
            \begin{itemize}
                \item \textbf{Responsibility:} The "engine" that handles the core logic.
                \item \textbf{Action:} Takes Celsius as an argument, performs the calculation, and returns the result.
            \end{itemize}
        \item \texttt{display\_result()}:
            \begin{itemize}
                \item \textbf{Responsibility:} Manages the final output formatting.
                \item \textbf{Action:} Takes two arguments and prints a formatted string to the console.
            \end{itemize}
    \end{itemize}
    \end{block}
\end{frame}

\begin{frame}[fragile]
    \frametitle{Putting It All Together: Execution Flow \& Takeaways}
    
    \begin{block}{The "Main" Script: The Orchestrator}
    The code outside the functions runs sequentially, calling each function to perform its task in the correct order.
    
    \begin{enumerate}
        \item \texttt{celsius = get\_celsius\_input()} \\
        Calls the input function and stores its \texttt{return} value.
        
        \item \texttt{fahrenheit = convert\_c\_to\_f(celsius)} \\
        Passes the stored Celsius value to the conversion function and stores the new result.
        
        \item \texttt{display\_result(celsius, fahrenheit)} \\
        Passes both values to the display function to print the final output.
    \end{enumerate}
    \end{block}
    
    \bigskip
    \textbf{Key Takeaways:}
    \begin{itemize}
        \item \textbf{Clarity:} The main script reads like a high-level summary of the program's steps.
        \item \textbf{Modularity:} To change the conversion formula (e.g., to Kelvin), you only need to edit \texttt{convert\_c\_to\_f}.
        \item \textbf{Reusability:} The \texttt{convert\_c\_to\_f} function is a self-contained tool that can be used in other programs.
    \end{itemize}
\end{frame}

\begin{frame}[fragile]
    \frametitle{Chapter 7: Summary - Core Mechanics}
    
    \begin{block}{1. Defining vs. Calling: The Two-Step Process}
        A function is a block of code that only runs when you explicitly tell it to.
        \begin{itemize}
            \item \textbf{Definition (\texttt{def})}: This is the blueprint for your function. The code inside is \textit{not} executed yet.
            \item \textbf{Calling}: This is when you execute the function's code by its name, followed by parentheses \texttt{()}.
        \end{itemize}
    \end{block}
    
    \begin{block}{2. Parameters: The Function's "In-box"}
        Parameters are placeholders that make functions flexible and reusable. You provide actual values (arguments) when you call the function.
        \begin{lstlisting}[language=Python]
% 'user_name' is a parameter (placeholder)
def say_hello(user_name):
    print(f"Hello, {user_name}!")

say_hello("Charlie") % "Charlie" is an argument
say_hello("Dana")    % "Dana" is an argument
        \end{lstlisting}
    \end{block}
    
\end{frame}

\begin{frame}[fragile]
    \frametitle{Chapter 7: Summary - Return Values \& Scope}
    
    \begin{block}{3. The \texttt{return} Statement: Getting a Result Back}
        The \texttt{return} keyword sends a value back from the function to the code that called it. This allows you to use the result elsewhere.
        \begin{lstlisting}[language=Python]
def add(a, b):
    result = a + b
    return result % Sends the value of 'result' back

sum_val = add(5, 3) % The returned value is stored
print(f"The result is: {sum_val}")
        \end{lstlisting}
    \end{block}
    
    \begin{block}{4. Local Scope: What Happens in a Function, Stays in a Function}
        Variables created inside a function are \textbf{local}. They only exist within that function and are destroyed when it finishes.
        \begin{lstlisting}[language=Python]
def calculate_area(width, height):
    area = width * height % 'area' is a local variable
    return area

rectangle_area = calculate_area(10, 5)
% print(area) % This would cause a NameError
        \end{lstlisting}
    \end{block}
\end{frame}

\begin{frame}[fragile]
    \frametitle{Chapter 7: Summary - The Big Picture}
    
    \begin{block}{5. Decomposition: Thinking Like a Programmer}
        The practice of breaking down a large, complex problem into smaller, more manageable sub-problems. Each sub-problem is then solved with a dedicated function.
        \begin{itemize}
            \item This makes code easier to read, test, and debug.
            \item Example: The temperature converter was broken into \texttt{get\_celsius\_input()}, \texttt{convert\_c\_to\_f()}, and \texttt{display\_result()}.
        \end{itemize}
    \end{block}
    
    \begin{alertblock}{Key Takeaways}
        \begin{itemize}
            \item Functions are defined with \textbf{\texttt{def}} and executed when \textbf{called}.
            \item \textbf{Parameters} are placeholders for inputs, making functions reusable.
            \item The \textbf{\texttt{return}} keyword sends a value back from a function.
            \item Variables inside a function are \textbf{local} to that function's scope.
            \item \textbf{Decomposition} helps manage complexity by breaking tasks into smaller functions.
        \end{itemize}
    \end{alertblock}
\end{frame}


\end{document}